% !TEX program = pdflatex
% Carte mentale - Les suites (Première)
% Compilation : pdflatex carte-mentale-suites.tex
\documentclass[11pt]{article}

\usepackage[utf8]{inputenc}
\usepackage[T1]{fontenc}
\usepackage[french]{babel}
\usepackage[a4paper,landscape,margin=1.1cm]{geometry}
\usepackage{lmodern}
\usepackage{amsmath,amssymb}
\usepackage{xcolor}
\usepackage{tikz}
\usepackage{bm}
\usetikzlibrary{arrows.meta,fit,backgrounds}

%------------------------------------------------------------
% Couleurs
%------------------------------------------------------------
\definecolor{cCentre}{RGB}{30,80,140}
% Bulles
\definecolor{c1}{RGB}{200,70,60}
\definecolor{c2}{RGB}{210,130,40}
\definecolor{c3}{RGB}{40,150,110}
\definecolor{c4}{RGB}{70,110,200}
\definecolor{c5}{RGB}{150,90,170}
\definecolor{c6}{RGB}{200,90,130}
% Grands blocs
\definecolor{cQ}{RGB}{120,100,60}
\definecolor{cA}{RGB}{55,90,135}
\definecolor{cG}{RGB}{60,140,90}

\pagestyle{empty}

%------------------------------------------------------------
\begin{document}

\begin{center}
  {\color{cCentre}\LARGE\bfseries Carte mentale -- Les suites}\\[3pt]
  {\itshape Première -- synthèse de ce qui a été vu}
\end{center}

\vspace{0.1cm}

\begin{center}
\begin{tikzpicture}[
  every node/.style={align=left},
  branch/.style={rounded corners=5pt, draw=#1, line width=0.9pt,
                 fill=#1!8, text width=4.1cm, inner sep=6pt, font=\small},
  banner/.style={rounded corners=5pt, fill=#1, text=white,
                 font=\bfseries, inner sep=5pt, align=center},
  panel/.style={rounded corners=11pt, draw=#1, line width=1.4pt,
                fill=#1!4, inner sep=11pt},
  link/.style={line width=1.3pt}
]

%====================== NOEUD CENTRAL ======================
\node[fill=cCentre, text=white, rounded corners=12pt, inner sep=12pt,
      font=\large\bfseries, align=center] (C) at (0,2.2) {LES\\ SUITES};

%====================== BLOC « QUELCONQUE » (gauche) ======================
\node[branch=c1, text width=6.4cm] (b1) at (-8.6,4.2) {%
  \begin{center}
    $u_n$ et $u(n)$ veulent dire la même chose.
  \end{center}  

  \vspace{5pt}
  {\color{c1}\bfseries Deux façons de définir}\\[2pt]
  $\bullet$ \textbf{Explicite} : 
  \\\underline{exemple :} $u_n = 6n + 2$\\
  \textbf{$\bm{u_n}$ est définie en fonction de $\bm{n}$.}

  \vspace{5pt}
  $\bullet$ \textbf{Récurrence} : $u_0$ donné,\\
  \phantom{$\bullet$ }\underline{exemple :} $u_{n+1} = \dfrac{1}{2}u_n - 2$};

\node[branch=c2, text width=6.4cm] (b2) at (-8.6,0.3) {%
  {\color{c2}\bfseries Calculer un terme}\\[3pt]
  $\bullet$ \textbf{Explicite} : remplacer $n$ par sa valeur\\
  $\bullet$ \textbf{Récurrence} : de proche en proche $u_0, u_1, u_2, \dots$};

%====================== BLOC « ARITHMÉTIQUE » (droite) ======================
\node[branch=c3] (b3) at (5.6,4.2) {%
  {\color{c3}\bfseries Définition}\\[3pt]
  On \textbf{ajoute} toujours le même nombre $r$ (\textbf{raison}) au terme précédent :\\[2pt]
  \hspace*{0.3cm}$u_{n+1} = u_n + r$};

\node[branch=c4, text width=5cm] (b4) at (10.8,4.2) {%
  {\color{c4}\bfseries Formules}\\[3pt]
  $\bullet$ \textbf{Récurrence :} $u_{n+1} = u_n + r$\\

  \vspace{5pt}
  $\bullet$ \textbf{Explicite :} $u_n = u_0 + nr$\\

  \vspace{5pt}
  $\bullet$ \it{ou cas particulier} 
  \\$u_n = u_p + (n-p)\,r$};

\node[branch=c5, text width=5cm] (b5) at (5.6,0.3) {%
  {\color{c5}\bfseries Représentation}\\[3pt]
  $\bullet$ $n$ en abscisses, $u_n$ en ordonnées\\
  $\bullet$ Les points sont alignés et forment une \textbf{droite}.};

\node[branch=c6] (b6) at (11.3,0.3) {%
  {\color{c6}\bfseries Sens de variation}\\[3pt]
  Signe de la raison $r$ :\\
  $\bullet$ $r>0$ : \textbf{croissante}\\
  $\bullet$ $r<0$ : \textbf{décroissante}\\
  $\bullet$ $r=0$ : \textbf{constante}};

%====================== BLOC « GÉOMÉTRIQUE » (bas, à compléter) ======================
\node[panel=cG, dashed, fill=cG!4, text width=6cm, align=center,
      font=\normalsize] (pG) at (0,-4.1) {%
  \vspace{2pt}{\Huge\color{cG!85!black}\bfseries ?}\\[6pt]
  \itshape (dernier bloc à compléter)\\[2pt]};

%====================== PANNEAUX + BANNIÈRES (arrière-plan) ======================
\begin{scope}[on background layer]
  \node[panel=cQ, fit=(b1)(b2)] (pQ) {};
  \node[panel=cA, fit=(b3)(b4)(b5)(b6)] (pA) {};

  % Liens depuis le centre
  \draw[link, cQ] (C) -- (pQ.east);
  \draw[link, cA] (C) -- (pA.west);
  \draw[link, cG] (C) -- (pG.north);
\end{scope}

% Bannières de titre
\node[banner=cQ, anchor=south] at ([yshift=4pt]pQ.north)
     {Suites quelconques\\[1pt]{\footnotesize\mdseries vrai pour \emph{toute} suite}};
\node[banner=cA, anchor=south] at ([yshift=4pt]pA.north)
     {Suites arithmétiques};
\node[banner=cG, anchor=south] at ([yshift=4pt]pG.north)
     {Suites géométriques};

\end{tikzpicture}
\end{center}

\end{document}
